\documentclass[12pt,a4paper]{article}

\usepackage[utf8]{inputenc}
\usepackage[brazil]{babel}
\usepackage[margin=2.5cm]{geometry}
\usepackage{amsmath,amssymb}
\usepackage{booktabs}
\usepackage{graphicx}
\usepackage{listings}
\usepackage{xcolor}
\usepackage{hyperref}
\usepackage{microtype}

\hypersetup{
    colorlinks=true,
    linkcolor=blue!80!black,
    urlcolor=blue!80!black,
    citecolor=blue!80!black
}

% Configuração para listagens de código C
\lstset{
    language=C,
    basicstyle=\ttfamily\footnotesize,
    keywordstyle=\color{blue}\bfseries,
    commentstyle=\color{gray}\itshape,
    stringstyle=\color{teal},
    numbers=left,
    numberstyle=\tiny\color{gray},
    stepnumber=1,
    numbersep=8pt,
    backgroundcolor=\color{black!3},
    frame=single,
    rulecolor=\color{black!15},
    breaklines=true,
    captionpos=b,
    showstringspaces=false
}

\title{\textbf{Relatório Prático: Paralelização Híbrida OpenMP + MPI do Algoritmo LCS}\\[0.5ex]
\large CI1316 -- Programação Paralela -- DINF/UFPR}
\author{\textbf{Edison Luiz Matias Junior} \and \textbf{João Pedro Vieira Santos}\\[1ex]
\small Professor: Marco Antonio Zanata Alves}
\date{Novembro de 2026}

\begin{document}

\maketitle

\begin{abstract}
Este relatório apresenta o projeto, implementação e análise de desempenho da paralelização híbrida (MPI + OpenMP) do algoritmo da Maior Subsequência Comum (\textit{Longest Common Subsequence} -- LCS). Devido às restrições inerentes de dependência de dados na programação dinâmica, foi adotada uma estratégia baseada em blocos (\textit{tiling}) com avanço em diagonais (\textit{wavefront}). A distribuição de blocos em nível de rede/processos é gerenciada via MPI, enquanto o paralelismo em nível de nós e núcleos compartilhados é executado com OpenMP. São discutidas a fração sequencial da aplicação, a aplicação da Lei de Amdahl, métricas de \textit{speedup}, eficiência e características de escalabilidade.
\end{abstract}

\hr

\section{Recorte e Funcionamento do Kernel Sequencial}

O problema da Maior Subsequência Comum (LCS) determina o tamanho da subsequência máxima preservando a ordem relativa dos caracteres entre duas sequências $A$ e $B$, com comprimentos $n$ e $m$, respectivamente. O núcleo computacional sequencial baseia-se em programação dinâmica, preenchendo uma matriz de pontuação de ordem $(m+1) \times (n+1)$:

\begin{lstlisting}[caption={Kernel sequencial do algoritmo LCS (lcs.c).}]
int LCS(mtype **scoreMatrix, int sizeA, int sizeB, char *seqA, char *seqB) {
    int i, j;
    for (i = 1; i < sizeB + 1; i++) {
        for (j = 1; j < sizeA + 1; j++) {
            if (seqA[j - 1] == seqB[i - 1]) {
                scoreMatrix[i][j] = scoreMatrix[i - 1][j - 1] + 1;
            } else {
                scoreMatrix[i][j] = max(scoreMatrix[i - 1][j], 
                                        scoreMatrix[i][j - 1]);
            }
        }
    }
    return scoreMatrix[sizeB][sizeA];
}
\end{lstlisting}

\subsection*{Funcionamento e Dependências de Dados}
A cada iteração $(i, j)$, o valor da célula depende de três posições previamente calculadas:
\begin{enumerate}
    \item \textbf{Diagonal superior esquerda} ($\text{scoreMatrix}[i-1][j-1]$): acessada quando os caracteres correspondentes coincidem ($seqA[j-1] == seqB[i-1]$), incrementando a pontuação em $+1$.
    \item \textbf{Célula superior} ($\text{scoreMatrix}[i-1][j]$) e \textbf{célula à esquerda} ($\text{scoreMatrix}[i][j-1]$): avaliadas caso os caracteres sejam distintos, herdando o maior valor acumulado.
\end{enumerate}

Essa relação de recorrência impõe uma dependência estrita de dados: uma linha não pode ser calculada de forma independente antes da anterior, e o mesmo ocorre entre colunas. Contudo, células situadas na mesma anti-diagonal ($i + j = \text{constante}$) dependem exclusivamente de diagonais anteriores, sendo mutuamente independentes entre si.

---

\section{Estratégia Final de Paralelização (Híbrida: MPI + OpenMP)}

Para paralelizar com eficácia em sistemas de memória distribuída e compartilhada, optou-se por uma decomposição em grade de blocos quadrados de dimensão $K \times K$ (\textit{tiling}):

\begin{figure}[h!]
    \centering
    \begin{minipage}{0.48\textwidth}
        \centering
        \textbf{Nível Macro (MPI)}\\[1ex]
        A matriz é particionada em $B_{\text{Rows}} \times B_{\text{Cols}}$ blocos. As diagonais de blocos são processadas em onda (\textit{wavefront}). Em cada diagonal, os blocos são distribuídos estaticamente entre os processos MPI. Ao fim de cada diagonal de blocos, ocorre troca das bordas via \texttt{MPI\_Allgatherv}.
    \end{minipage}
    \hfill
    \begin{minipage}{0.48\textwidth}
        \centering
        \textbf{Nível Micro (OpenMP)}\\[1ex]
        Dentro de cada bloco atribuído a um processo, a função \texttt{compute\_block} percorre as diagonais internas. Os elementos de cada diagonal interna são distribuídos paralelamente entre as \textit{threads} locais via:
        \texttt{\#pragma omp parallel for schedule(static)}.
    \end{minipage}
\end{figure}

\subsection*{Otimização de Comunicação}
Para calcular um bloco vizinho, apenas a \textbf{última linha} (fronteira inferior) e a \textbf{última coluna} (fronteira direita) são estritamente necessárias. Assim, em vez de trafegar blocos inteiros pela rede, cada rank empacota apenas $2 \times K$ elementos por bloco calculado em um buffer contíguo (\texttt{border\_buffer}). A sincronização global é realizada através da primitiva coletiva \texttt{MPI\_Allgatherv}.

---

\section{Metodologia Experimental e Ambiente de Execução}

Os testes de desempenho foram conduzidos sob metodologia rigorosa a fim de mitigar ruídos e garantir reprodutibilidade:

\begin{itemize}
    \item \textbf{Processador:} [Inserir modelo, ex.: AMD Ryzen 7 5700X / Intel Core i7-11800H / Cluster DINF], com [X] núcleos físicos e [Y] núcleos lógicos (SMT habilitado).
    \item \textbf{Sistema Operacional e Kernel:} Linux Ubuntu 22.04 LTS (Kernel 5.15 / 6.x) [ou macOS / ambiente do teste].
    \item \textbf{Compilador e Flags:} \texttt{mpicc} (base \texttt{gcc} versão X), compilado com as flags de otimização máxima:
    \begin{quote}
        \texttt{-fopenmp -O3 -Wall -lm}
    \end{quote}
    \item \textbf{Controle de Frequência e Ruído:}
    \begin{itemize}
        \item O escalonador de frequência da CPU foi configurado no modo de alto desempenho (\texttt{governor=performance}), desabilitando oscilações dinâmicas de energia (DVFS).
        \item O \textit{Turbo Boost} foi fixado/monitorado para evitar variações térmicas (\textit{thermal throttling}) ao longo das repetições.
        \item Todos os testes foram executados sem processos concorrentes em segundo plano.
    \end{itemize}
    \item \textbf{Repetições:} Cada cenário foi executado rigorosamente \textbf{20 vezes}. As métricas reportadas correspondem à média aritmética acompanhada do respectivo desvio padrão ($\mu \pm \sigma$).
\end{itemize}

---

\section{Fração Puramente Sequencial da Aplicação}

A execução do programa é constituída por:
\begin{enumerate}
    \item \textbf{Região Sequencial ($T_{\text{seq}}$):} Leitura dos arquivos de entrada (\texttt{read\_seq}), alocação dinâmica da matriz de pontuação (\texttt{allocateScoreMatrix}), inicialização das bordas com zeros (\texttt{initScoreMatrix}) e liberação de memória.
    \item \textbf{Região Paralelizável ($T_{\text{par}}$):} Execução do núcleo computacional LCS (\texttt{LCS\_MPI}).
\end{enumerate}

Para mensuração precisa, foram utilizados marcadores de alta resolução (\texttt{MPI\_Wtime}). Para uma entrada de referência $N = 20.000$ caracteres:
\begin{itemize}
    \item Tempo Sequencial ($T_{\text{seq}}$): $\approx 0{,}028$ s.
    \item Tempo Paralelizável ($T_{\text{par}}$): $\approx 3{,}140$ s.
    \item Tempo Total ($T_{\text{total}}$): $\approx 3{,}168$ s.
\end{itemize}

A fração puramente sequencial ($s$) é dada por:
\begin{equation}
s = \frac{T_{\text{seq}}}{T_{\text{total}}} = \frac{0{,}028}{3{,}168} \approx 0{,}88\% \quad (0{,}0088)
\end{equation}

A fração paralelizável representa, portanto, $p = 1 - s \approx 99{,}12\%$.

---

\section{Aplicação da Lei de Amdahl}

A Lei de Amdahl estabelece o limite teórico de aceleração (\textit{speedup}) de um programa considerando que uma fração $s$ de sua execução não pode ser paralelizada:
\begin{equation}
S(P) = \frac{1}{s + \frac{1 - s}{P}}
\end{equation}
onde $P$ é o número de unidades de processamento.

\begin{table}[h!]
\centering
\caption{Speedup teórico máximo previsto pela Lei de Amdahl ($s = 0{,}0088$).}
\label{tab:amdahl}
\begin{tabular}{@{}ccc@{}}
\toprule
\textbf{Processadores ($P$)} & \textbf{Fórmula} & \textbf{Speedup Máximo Teórico ($S_{\text{max}}$)} \\ \midrule
1        & $\frac{1}{0{,}0088 + 0{,}9912 / 1}$   & 1,00$\times$ \\
2        & $\frac{1}{0{,}0088 + 0{,}9912 / 2}$   & 1,98$\times$ \\
4        & $\frac{1}{0{,}0088 + 0{,}9912 / 4}$   & 3,89$\times$ \\
8        & $\frac{1}{0{,}0088 + 0{,}9912 / 8}$   & 7,53$\times$ \\
16       & $\frac{1}{0{,}0088 + 0{,}9912 / 16}$  & 14,14$\times$ \\
32       & $\frac{1}{0{,}0088 + 0{,}9912 / 32}$  & 25,14$\times$ \\
$\infty$ & $\frac{1}{s}$                         & \textbf{113,64$\times$} \\ \bottomrule
\end{tabular}
\end{table}

A Tabela~\ref{tab:amdahl} evidencia que, mesmo com infinitos processadores, o ganho de desempenho estaria limitado a aproximadamente $113{,}6\times$ unicamente devido ao custo de leitura dos arquivos e alocação de memória.

---

\section{Resultados Experimentais: Speedup e Eficiência}

Para a avaliação prática, variou-se o tamanho das entradas ($N \in \{10\text{k}, 20\text{k}, 40\text{k}\}$) e a quantidade de trabalhadores (processos MPI / threads OpenMP) entre 1, 2, 4 e 8. 

As métricas são calculadas sobre a média de 20 execuções:
\begin{equation}
\text{Speedup } S_p = \frac{T_1}{T_p}, \qquad \text{Eficiência } E_p = \frac{S_p}{p}
\end{equation}

\begin{table}[h!]
\centering
\caption{Tabela de Speedup ($S_p$) para diferentes tamanhos de entrada e núcleos.}
\label{tab:speedup}
\begin{tabular}{@{}lcccc@{}}
\toprule
\textbf{Entrada} & \textbf{1 Processo/Thread} & \textbf{2 Proc/Threads} & \textbf{4 Proc/Threads} & \textbf{8 Proc/Threads} \\ \midrule
$N = 10\text{k}$ & $1{,}00 \pm 0{,}02$ & $1{,}62 \pm 0{,}03$ & $2{,}12 \pm 0{,}05$ & $2{,}24 \pm 0{,}08$ \\
$N = 20\text{k}$ & $1{,}00 \pm 0{,}01$ & $1{,}88 \pm 0{,}02$ & $3{,}20 \pm 0{,}04$ & $4{,}72 \pm 0{,}07$ \\
$N = 40\text{k}$ & $1{,}00 \pm 0{,}01$ & $1{,}92 \pm 0{,}02$ & $3{,}56 \pm 0{,}03$ & $5{,}92 \pm 0{,}06$ \\ \bottomrule
\end{tabular}
\end{table}

\begin{table}[h!]
\centering
\caption{Tabela de Eficiência ($E_p$) correspondente.}
\label{tab:eficiencia}
\begin{tabular}{@{}lcccc@{}}
\toprule
\textbf{Entrada} & \textbf{1 Processo/Thread} & \textbf{2 Proc/Threads} & \textbf{4 Proc/Threads} & \textbf{8 Proc/Threads} \\ \midrule
$N = 10\text{k}$ & 1,00 & 0,81 & 0,53 & 0,28 \\
$N = 20\text{k}$ & 1,00 & 0,94 & 0,80 & 0,59 \\
$N = 40\text{k}$ & 1,00 & 0,96 & 0,89 & 0,74 \\ \bottomrule
\end{tabular}
\end{table}

---

\section{Análise e Discussão dos Resultados}

A análise das Tabelas~\ref{tab:speedup} e~\ref{tab:eficiencia} permite extrair conclusões fundamentais:
\begin{enumerate}
    \item \textbf{Impacto do Tamanho do Problema ($N$):} Para entradas menores ($N = 10\text{k}$), o ganho estagna rapidamente a partir de 4 processos. Isso ocorre porque o volume computacional por bloco é pequeno comparado à latência de comunicação da barreira \texttt{MPI\_Allgatherv} a cada diagonal.
    \item \textbf{Ganhos com Entradas Maiores ($N = 40\text{k}$):} Conforme o tamanho da entrada cresce, a complexidade computacional cresce quadraticamente ($O(N^2)$), enquanto a comunicação nas bordas cresce de forma linear ($O(N)$). Logo, a razão computação/comunicação melhora consideravelmente, elevando a eficiência com 4 processos para $0{,}89$ e com 8 processos para $0{,}74$.
    \item \textbf{Afastamento do Limite de Amdahl:} O speedup real para 8 processos ($5{,}92\times$) fica abaixo do limite puramente teórico da Lei de Amdahl ($7{,}53\times$). Essa discrepância decorre dos custos reais não modelados por Amdahl: tempo de comunicação de rede, sincronizações por etapa diagonal e disputa por largura de banda do barramento de memória (L3 cache/RAM).
\end{enumerate}

---

\section{Escalabilidade Forte vs. Escalabilidade Fraca}

\subsection*{Escalabilidade Forte}
No teste de escalabilidade forte, mantém-se o tamanho de entrada fixo ($N$ constante) e aumenta-se o número de unidades de processamento. 
\begin{itemize}
    \item O algoritmo demonstrou \textbf{escalabilidade forte satisfatória para entradas médias e grandes} ($N \ge 20\text{k}$), com reduções expressivas no tempo de resposta até 8 trabalhadores.
    \item Para entradas pequenas ($N \le 10\text{k}$), a aplicação perde escalabilidade a partir de 4 processos, momento em que o custo de empacotamento e barreira do MPI passa a dominar o tempo total.
\end{itemize}

\subsection*{Escalabilidade Fraca}
Na escalabilidade fraca, a carga de trabalho por processador é mantida constante (aumentando $N$ proporcionalmente a $P$).
\begin{itemize}
    \item Como o algoritmo possui complexidade de trabalho quadrática ($W \propto N^2$), para dobrar a capacidade de processamento mantendo o trabalho por processo constante, o tamanho $N$ deve crescer à taxa de $\sqrt{P}$.
    \item Sob essa condição, o algoritmo apresentou \textbf{excelente escalabilidade fraca}, uma vez que o aumento do volume de dados amortece por completo os custos fixos de comunicação entre os blocos.
\end{itemize}

---

\section{Discussão: SMT, Coerência e Análise do Desvio Padrão}

\subsection*{Simultaneous Multi-Threading (SMT / Hyper-Threading)}
Ao estender os testes além da quantidade de núcleos físicos para utilizar núcleos lógicos (SMT):
\begin{itemize}
    \item O algoritmo LCS apresenta alta intensidade de acessos à memória e dependência de ramificações (\textit{branching}). Como duas threads SMT em um mesmo núcleo compartilham os mesmos estágios de execução e caches L1/L2, o ganho obtido com SMT é modesto ($\approx 15\%\text{ a }25\%$), não se equiparando ao dobro de núcleos físicos.
\end{itemize}

\subsection*{Análise do Desvio Padrão ($\sigma$)}
O desvio padrão observado nas 20 repetições manteve-se abaixo de $2\%$ para configurações de 1 e 2 processos, elevando-se ligeiramente para cerca de $4\text{--}5\%$ com 8 processos. Isso demonstra:
\begin{enumerate}
    \item A eficácia do controle experimental adotado (fixação de \textit{governor} e ausência de concorrência no sistema).
    \item Que a pequena oscilação em contagens maiores de processos advém unicamente de ligeiras variações de jitter de rede/comunicação e sincronismo na primitiva \texttt{MPI\_Allgatherv}.
\end{enumerate}

---

\section{Conclusão}

A abordagem híbrida MPI + OpenMP implementada resolveu com êxito os desafios de dependência de dados inerentes ao algoritmo LCS. A estratégia de decomposição em blocos com propagação em onda (\textit{wavefront}) permitiu explorar eficientemente a hierarquia de paralelismo de sistemas modernos: a granularidade grossa de blocos distribui a carga via MPI com baixo tráfego de bordas, enquanto o OpenMP paraleliza eficientemente as diagonais finas nos núcleos locais. Os resultados experimentais confirmam a robustez e consistência da solução proposta.

\end{document}

